%2025 - 2b
\chapter{Variétés différentielles}
\section{02b Chapitre 2 : Variétés différentielles}

Spontanément on pense à un espace comme étant euclidien comme ayant la géométrie euclidienne qui est la géométrie naturelle. Il y a une particularité dans ce type d'espace, c'est la structure affine : est associé à cet espace un espace vectorielle

espace : ensemble de point \qquad $\to$ espace vectorielle
\[
(A,B)\ \to\ \vec{AB}
\]
Un tel espace est très particulier. Tout dit que notre espace temps n'a pas une telle structure. Beaucoup de chose en physique repose sur l'existence de vecteur d'espace.

l'évolution des systèmes physiques est continue et différentiable. Lorsqu'on laisse tomber un objet matériel, il passe continument par tout les points de sa trajectoire.

%Cas des trajectoires des corps physiques

%Liberté de choix d'un système de coordonnées

%Nécessité d'une notion {\bf intrinsèque} de continuité et de différentiabilité

%Espaces topologiques :

%-notion de topologie (ensembles "ouverts"), de voisinage, de limite, de continuité


